\backmatter

% INSA Stuff
\newpage %last page must be even
\pagestyle{empty} % No headings in folio
% \sffamily
\newgeometry{margin=2.5cm}
\ULforem

\newgeometry{top=.8cm,bottom=1cm,right=1.5cm,left=1.5cm}

\begin{center}
\includegraphics[height=1.9cm]{INSA.png}
\end{center}

\vspace{.3cm}
\footnotesize
\centering
\begin{tabular}{|p{8.6cm}p{8.6cm}|}	
	\multicolumn{2}{p{17.2cm}}{\center{\normalsize{\underline{FOLIO ADMINISTRATIF}}}}\\[-.5cm]
	\multicolumn{2}{p{17.2cm}}{\center{\normalsize{\dotuline{THESE DE L’INSA LYON, MEMBRE DE L’UNIVERSITE DE LYON}}}}\\
	\multicolumn{2}{p{17.2cm}}{}\\
	\hline
	&\\
	NOM : \textbf{RUSHTON}&DATE de SOUTENANCE : \textbf{23/02/2027}\\
	&\\
	Prénoms : \textbf{Thomas Albert}&\\
	&\\
	\multicolumn{2}{|p{17.2cm}|}{TITRE : \textbf{A
	Sample-Synchronous Distributed Audio System for Spatial and
	Immersive Applications}}\\
	&\\
	&\\
	&\\
	NATURE : \textbf{Doctorat}&Numéro d'ordre : \textbf{xxxxxxx}\\
	&\\
	\multicolumn{2}{|p{17.2cm}|}{\'Ecole Doctorale : \textbf{Informatique et Mathématiques}}\\
	&\\
	Spécialité : \textbf{Informatique}&\\
	&\\
	&\\
	R\'ESUM\'E :&\\% attention, si tu me changes change moi aussi dans l'intro
	\multicolumn{2}{|p{17.2cm}|}{L'accès aux systèmes audio spatialisés à grande échelle demeure difficile, malgré l'intérêt considérable que leur portent la recherche comme l'industrie. Le coût des interfaces audio à très grand nombre de canaux et des systèmes de traitement numérique du signal puissants rend ces installations hors de portée de pratiquement tous les acteurs, à l'exception des institutions les mieux dotées. La dimension auditive de l'expérience immersive reste ainsi relativement peu explorée dans l'art sonore et la réalité virtuelle, en raison du seuil d'entrée élevé des techniques de synthèse de champ sonore que sont la synthèse de champ d'ondes (WFS) et l'ambisonie d'ordre supérieur (HOA).}\\
        \multicolumn{2}{|p{17.2cm}|}{Cette thèse décrit un système modulaire et extensible de spatialisation sonore distribuée, qui tire parti de l'une des nombreuses plateformes audio embarquées à bas coût apparues au cours de la dernière décennie environ. Un réseau ethernet de tels dispositifs fait office de processeur de signal numérique distribué et coordonné. Grâce à la récupération de l'horloge média et de l'instant de présentation, déduite d'échanges d'informations temporelles fondés sur le protocole PTP (Precision Time Protocol) de la norme IEEE 1588, les modules matériels restent synchronisés à l'échelle de la sub-microseconde. Le système repose en outre sur un serveur audio intégré à une station de travail audionumérique, ainsi que sur des protocoles conçus pour la transmission en réseau de flux audio, de données de contrôle et de la découverte des dispositifs, dans un contexte à contraintes temporelles strictes. Des algorithmes de spatialisation sonore, sous forme parallélisée, s'appuient sur ce socle synchronisé avec précision.}\\
	&\\
	\multicolumn{2}{|p{17.2cm}|}{MOTS-CL\'ES : \textbf{audio spatialisé, protocoles réseau, traitment du signal
distribué, systèmes embarqués, échange temporel}}\\
	&\\
	&\\
	\multicolumn{2}{|p{17.2cm}|}{
	\begin{tabular}{@{}ll}%@{} pour avoir aucune bordure
	Laboratoire(s) de recherche : &\textbf{CITI}\\
	\end{tabular}}\\
	&\\
  \multicolumn{2}{|p{17.2cm}|}{Directeur de thèse :
  \textbf{Tanguy Risset}}\\
	&\\
  \multicolumn{2}{|p{17.2cm}|}{Président du Jury :
  \textbf{Jean-Marie GORCE}}\\
	&\\
  \multicolumn{2}{|p{17.2cm}|}{Composition du Jury :}\\
  \multicolumn{2}{|p{17.2cm}|}{\hspace{1cm}\textbf{Stefania
  SERAFIN, Alain BONARDI, Cristina ROTTONDI, Jens AHRENS, Tanguy RISSET,
  Romain MICHON}}\\
	&\\
  \hline
\end{tabular}
